\documentclass[12pt]{article}
\usepackage{espd}
\docode{RU.МИФИ.52714-01 ТЗ 01-1}
% \makeapproval{Наименование}{Вид документа}{Шифр ЛУ}{Шифр}
\newcommand{\univ}{Национальный исследовательский ядерный университет «МИФИ»}
\newcommand{\course}{Дисциплина «Генеративный искусственный интеллект: модели и приложения»}
\newcommand{\progname}{Компьютерная игра «Dreich Day»}
\newcommand{\makeapproval}[3]{%
\thispagestyle{empty}
\begin{center}\singlespacing
\univ\\[2mm]\course
\end{center}
\vspace{10mm}
\noindent\begin{minipage}[t]{0.48\textwidth}\singlespacing
СОГЛАСОВАНО\\[2mm]Студент группы Б25-527\\[6mm]\rule{30mm}{0.4pt} А.\,И.~Лобанов\\[2mm]«\rule{8mm}{0.4pt}»\rule{30mm}{0.4pt} 2026 г.
\end{minipage}\hfill
\begin{minipage}[t]{0.48\textwidth}\singlespacing\raggedleft
УТВЕРЖДАЮ\\[2mm]Преподаватель дисциплины\\[6mm]\rule{30mm}{0.4pt} А.\,А.~Кузьмин\\[2mm]«\rule{8mm}{0.4pt}»\rule{30mm}{0.4pt} 2026 г.
\end{minipage}
\vfill
\begin{center}
{\bfseries\MakeUppercase{\progname}}\\[6mm]
{\bfseries\MakeUppercase{#1}}\\[4mm]
ЛИСТ УТВЕРЖДЕНИЯ\\[4mm]
#2
\end{center}
\vfill
\begin{flushright}\singlespacing
Исполнитель:\\студент группы Б25-527\\Лобанов Алексей Игоревич
\end{flushright}
\vspace{10mm}
\begin{center}Москва 2026\end{center}
\clearpage
\thispagestyle{empty}
\noindent УТВЕРЖДЁН\\#2\\
\vfill
\begin{center}
{\bfseries\MakeUppercase{\progname}}\\[6mm]
{\bfseries\MakeUppercase{#1}}\\[4mm]
#3\\[4mm]
Листов \pageref{LastPage}
\end{center}
\vfill
\begin{center}Москва 2026\end{center}
\clearpage
\setcounter{page}{2}}

\begin{document}
\makeapproval{Техническое задание}{RU.МИФИ.52714-01 ТЗ 01-1-ЛУ}{RU.МИФИ.52714-01 ТЗ 01-1}

\section*{\centering АННОТАЦИЯ}
Настоящий документ является техническим заданием на разработку компьютерной игры «Dreich Day» — уютной пиксельной 2D-игры о временной петле, выполняемой в рамках большого домашнего задания по дисциплине «Генеративный искусственный интеллект: модели и приложения».

Документ содержит основания для разработки, назначение программы, требования к её функциональным характеристикам, надёжности, условиям эксплуатации и составу технических средств, а также стадии и этапы разработки и порядок контроля и приёмки. Техническое задание подготовлено с применением генеративного ИИ в режиме мета-фасилитации: модель последовательно выполняла роли заказчика-геймера, архитектора и бизнес-аналитика, а исполнитель выступал владельцем продукта.

Документ оформлен в соответствии с ГОСТ 19.201-78 «ЕСПД. Техническое задание. Требования к содержанию и оформлению».
\clearpage
\tableofcontents
\clearpage

\section{Введение}
\subsection{Наименование программы}
Наименование программы — компьютерная игра «Dreich Day» (шотл. \textit{dreich} — «серый, сырой, унылый день»).
\subsection{Краткая характеристика области применения}
Программа относится к классу развлекательных браузерных игр жанра «симулятор жизни с элементами головоломки-дедукции». Игрок управляет синоптиком Оливером Оверкастом, застрявшим во временной петле в шотландской деревне Глен-Дрих: день повторяется, но его условия (погода, события, настроение жителей) каждый раз выпадают заново. Цель игрока — разгадать скрытые правила поведения жителей и подготовить «идеальный день», в который деревенская команда по шинти выиграет финальный матч.

\section{Основания для разработки}
Основанием для разработки является задание на большое домашнее задание (БДЗ) по дисциплине «Генеративный искусственный интеллект: модели и приложения», выданное преподавателем А.\,А.~Кузьминым. Наименование темы разработки — «Разработка компьютерной игры с применением генеративного ИИ».

\section{Назначение разработки}
\subsection{Функциональное назначение}
Программа предоставляет игроку игровой мир (деревню) с собственным циклом дня, случайно генерируемыми условиями и персонажами, поведение которых подчиняется скрытым детерминированным правилам. Игрок наблюдает, ведёт дневник, влияет на мир через прогноз погоды в газете и выполняет задания.
\subsection{Эксплуатационное назначение}
Программа предназначена для свободного использования в браузере широким кругом пользователей (14+) и служит демонстрацией освоения технологии вайбкодинга на платформе WebSim.

\section{Требования к программе}
\subsection{Требования к функциональным характеристикам}
\subsubsection{Цикл дня и временная петля}
\begin{enumerate}
\item Игровой день длится с 6:00 до 18:00; один игровой час соответствует 50 секундам реального времени.
\item В 6:00 каждой петли генерируются условия дня: погода утром и после полудня (солнце, дождь, туман, ветер), знак гадюки Натайр на утреннем обряде, событие дня (ярмарка, визит Ангуса или его отсутствие) и стартовое настроение жителей.
\item В 18:00 либо по решению игрока (кровать Оливера) петля сбрасывается: показывается сводка дня, состояние мира возвращается к началу, сохраняются только дневник, счётчик петель и прогресс заданий.
\end{enumerate}
\subsubsection{Скрытые правила жителей}
\begin{enumerate}
\item Каждый житель имеет базовое расписание и набор скрытых правил вида «если [условия] — то [место, действие, изменение настроения]».
\item Правила хранятся как данные, что позволяет добавлять жителей без переписывания движка.
\item При одинаковых условиях и прогнозе жители ведут себя одинаково (правила детерминированы).
\end{enumerate}
\subsubsection{Прогноз в газете}
Игрок до 9:00 один раз за петлю выбирает прогноз погоды (солнце, дождь, туман, ветер). Жители строят по нему планы; в 12:00 прогноз проверяется: при совпадении настроение жителей растёт, при несовпадении — падает.
\subsubsection{Дневник и наблюдения}
Наблюдения записываются автоматически; правило считается открытым после двух наблюдений в разных петлях. Неоткрытые правила отображаются как «???».
\subsubsection{Задания дня}
Каждую петлю в утренней газете публикуется одно выполнимое при текущих условиях задание с подсказкой; выполнение даёт подсказку к неоткрытому правилу и сохраняется между петлями.
\subsubsection{Финальный матч по шинти}
Кульминация — матч команды «Гадюки Глен-Дриха» (11 жителей и Оливер) против соседней деревни Силлерберн. Сила команды определяется настроением игроков, разрешёнными конфликтами, погодой и наличием подкупленного судьи. Победа возможна только при «идеальном дне».
\subsection{Требования к надёжности}
Программа не должна выдавать ошибок в консоли браузера при штатной работе; при невозможности загрузить спрайт используется процедурная отрисовка объекта.
\subsection{Условия эксплуатации}
Программа эксплуатируется в современном браузере с поддержкой HTML5 Canvas и JavaScript (ES2020). Специальной подготовки пользователя не требуется.
\subsection{Требования к составу и параметрам технических средств}
Персональный компьютер с клавиатурой, разрешение экрана не менее 1280×720, доступ к сети Интернет.
\subsection{Требования к информационной и программной совместимости}
Программа реализуется на HTML5, CSS и JavaScript без внешних библиотек и размещается на платформе WebSim. Прогресс сохраняется в localStorage браузера. Интерфейс и диалоги — на русском языке, имена собственные — шотландские.
\subsection{Специальные требования}
Графика — пиксельная, вид сверху, единая палитра «серого шотландского дня» с тёплым светом окон как акцентом. Спрайты персонажей, зданий, интерьеров и окружения создаются с применением генеративного ИИ.

\section{Требования к программной документации}
Состав документации: техническое задание (ГОСТ 19.201-78), пояснительная записка (ГОСТ 19.404-79). Документация оформляется по ГОСТ 19.106-78 в системе LaTeX.

\section{Технико-экономические показатели}
Разработка выполняется силами одного исполнителя с использованием бесплатных ресурсов платформы WebSim (ежедневные кредиты). Ориентировочная экономическая эффективность не рассчитывается.

\section{Стадии и этапы разработки}
\begin{longtable}{|p{10mm}|p{45mm}|p{95mm}|}\hline
№ & Этап & Содержание работ\\\hline\endhead
1 & ИИ-мозговой штурм & Генерация концепции игры в режиме мета-фасилитации\\\hline
2 & Техническое задание & Разработка настоящего ТЗ и промпта для WebSim\\\hline
3.1 & Итерация 1 «Петля» & Карта, цикл дня, погода, прогноз, три жителя, дневник\\\hline
3.2 & Итерация 2 «Цели» & Задания дня, спрайтовая графика, окружение, новые жители, визиты Ангуса\\\hline
3.3 & Итерация 3 «Вся деревня» & Остальные жители, истории и конфликты, разоблачение подкупа\\\hline
3.4 & Итерация 4 «Матч» & Мини-игра шинти, расчёт силы команды, балансировка\\\hline
3.5 & Итерация 5 «Финал» & Концовка, меню, звук, полировка\\\hline
4 & Документирование & Подготовка комплекта документации в LaTeX\\\hline
5 & Публикация & Размещение исходных кодов в системе контроля версий\\\hline
\end{longtable}

\section{Порядок контроля и приёмки}
Приёмка каждой итерации выполняется по критериям приёмки: отсутствие ошибок в консоли, разнообразие условий не менее чем в трёх из пяти подряд петель, детерминированность правил, возможность открыть каждое правило в дневнике, заметное влияние прогноза на поведение жителей, сохранение дневника после сброса петли. Итоговая приёмка — защита БДЗ перед преподавателем.

\changelist
\end{document}
